\section{Agenda de Pesquisa Prioritária para a Próxima Década}

A materialização desse horizonte depende da orquestração de seis eixos prioritários de Pesquisa e Desenvolvimento (P\&D), delineados para cobrir as lacunas expostas ao longo da obra:

\begin{enumerate}
    \item \textbf{Validação Clínica Multicêntrica:} Condução inadiável de Ensaios Clínicos Randomizados (ECRs) de longo prazo (mínimo de 5 anos) para atestar a taxa de sobrevida de reabilitações guiadas por gêmeos digitais \cite{duggal2026exploring}.
    \item \textbf{Padronização e Semântica (Open-Source):} Desenvolvimento de um formato de arquivo unificado e código aberto para gêmeos digitais odontológicos, integrando malhas STL, arquivos DICOM e metadados biomecânicos em um único contêiner criptografado.
    \item \textbf{Quantificação da Incerteza (\textit{Uncertainty Quantification}):} Adoção mandatória de simulações estocásticas (ex: expansão em caos polinomial) para modelar a variabilidade biológica inerente aos tecidos periodontais e ósseos.
    \item \textbf{Auditoria Algorítmica e \textit{Debiasing}:} Criação de consórcios de dados (Data Trusts) representativos da miscigenação craniofacial brasileira, prevenindo que IAs exógenas induzam ao erro sistemático em populações vulneráveis \cite{springer2026realising}.
    \item \textbf{Avaliação Tecnológica em Saúde (ATS):} Financiamento estatal para estudos de custo-efetividade que mensurem o impacto financeiro e orçamentário da incorporação dos gêmeos digitais na tabela do SUS.
\end{enumerate}
